% !TeX root = ../main.tex
\chapter{Complexes and the inverse of a gauge equation}\label{f:gauge-complex}
The scalar wave operator acquired its retarded inverse by imposing zero incoming data. In electromagnetism the same prescription does not distinguish potentials differing by the derivative of a gauge parameter of compact support. This is not a technical failure of an estimate: the operator is required by symmetry to kill a whole family of fields. What can replace an inverse? We first answer a question of linear algebra, for which neither a metric nor causality is needed.

\section{Cohomology and homotopy}
A graded vector space is a family of spaces $C^k$ indexed here by integers. A cochain complex is such a family with maps $D_k:C^k\to C^{k+1}$ satisfying $D_{k+1}D_k=0$. Write simply $D$ when the degree of the argument is known. A cocycle is a vector in $\ker D_k$; a coboundary is a vector in $\operatorname{im}D_{k-1}$. The zero composition is exactly what makes the quotient $H^k(C)=\ker D_k/\operatorname{im}D_{k-1}$ well defined. It records cocycles which cannot be removed by changes coming from one degree below. For instance $d^2=0$ makes differential forms a cochain complex. The angular one-form of Chapter~\ref{f:forms} represents a nonzero class, as its period detects. There is no requirement that a complex be exact: exactness at degree $k$ is the \emph{additional conclusion} $H^k=0$.

A cochain map $T:C\to C'$ consists of degree-preserving maps commuting with the differentials, so it maps closed vectors to closed vectors and boundaries to boundaries. If two cochain maps satisfy $T-S=D' h+hD$ with maps $h:C^k\to C'^{k-1}$, they are called chain homotopic. On a closed vector $v$ the difference is the boundary $D'hv$, so the induced maps on cohomology coincide. In particular if
\begin{equation}
 Dh+hD=1,\label{f:contraction}
\end{equation}
then every closed vector is exact: $v=D(hv)$ for $Dv=0$. The map $h$ is a contracting homotopy. With $1-\Pi$ on the right, where $\Pi^2=\Pi$ commutes with $D$, the same reasoning says every class has a representative in the retained part $\Pi C$. It does not authorize one to forget that part. A two-term complex $C^0\xrightarrow{P}C^1$ with no other nonzero terms has \eqref{f:contraction} exactly when $h_1$ is a two-sided inverse of $P$. Moving to a longer complex lets neighboring maps absorb the parts a single inverse cannot handle.

\section{What symmetry implies about an equation}
Let $K$ send infinitesimal proper gauge parameters to field variations at a solution, $J$ denote the linearized Euler operator, and $K^*$ the adjoint relative to the field--source and parameter--parameter-density pairings. Gauge-related fields are solutions together, so $JK=0$. Invariance of the action for compactly supported local gauge parameters gives the Noether identity $K^*\mathcal E=0$ (integrate the derivatives of the parameter by parts and use its arbitrariness). Linearizing at $\mathcal E=0$ gives $K^*J=0$. If a boundary is present, these statements require that the transformations preserve the boundary conditions and that the adjoint pairing leave no unwanted surface term. Nonzero endpoint or surface charges are not being declared zero by this argument.

For an irreducible linearized theory the resulting portion of the deformation complex is
\begin{equation}
 U\ \xrightarrow[{-1\ \to\ 0}]{K}\ V\
 \xrightarrow[{0\ \to\ 1}]{J}\ V^\vee\
 \xrightarrow[{1\ \to\ 2}]{K^*}\ U^\vee .\label{f:deformation-complex}
\end{equation}
The symbols now denote section spaces on stated support and boundary domains, with dual \emph{densities} for sources, not algebraic duals containing arbitrary functionals. At degree zero $\ker J/\operatorname{im}K$ is the space of linearized solutions up to proper gauge; at degree one a closed source $f$ obeys the compatibility condition $K^*f=0$. Each differential must map its claimed domain into the next. For example, $K\lambda=d\lambda$ in Maxwell theory. If the tangential potential is fixed on a wall, $j^*d\lambda=d(j^*\lambda)$ must vanish there; a parameter with zero trace suffices. Permission to change a frame in the bulk does not settle the wall domain.

This grading is particularly natural in the Batalin--Vilkovisky description \cite{CMR,Khavkine}. An ordinary gauge parameter is encoded by an odd variable called a ghost (an odd variable anticommutes with other odd variables). Fields and ghosts also have dual variables called antifields; in the usual function grading a field of degree $k$ has antifield of degree $-1-k$. An odd differential $Q$ has $Q^2=0$ and acts by the graded product rule $Q(ab)=(Qa)b+(-1)^{|a|}a\,Qb$. At a zero of $Q$ its derivative is a square-zero map on the tangent complex. The degrees indicated on \eqref{f:deformation-complex} are those of \emph{deformations}, with parameters one degree below field tangents, not degrees being assigned to ordinary commuting sections as if they were ghost coordinate functions. For a simple model the gauge part of the differential on functions sends an abelian potential $A$ to $dc$ and its ghost $c$ to zero. Thus it squares to zero on $A$ and leaves $dA$ invariant. Equations of motion require additional components, represented by antifields in the BV formulation. This describes what the linear complex measures without assuming existence of a full quantum BV theory with boundary. In a differential graded description of Ward identities, an exact local violation can be removed by changing a local functional; a violation in nonzero cohomology is a different assertion.

\section{Gauge fixing constructs a homotopy}
Choose a gauge condition $C:V\to U$ and a symmetric invertible weight $B:U\to U^\vee$. Form $P=J+C^*BC$, $R=CK$, and $R^*=K^*C^*$. Assume $J$ is formally symmetric on the paired domains. A source for the field can be solved using $P$, but only if the gauge condition and parameter equation can be solved as well. Fix either a retarded or an advanced support problem. Suppose the indicated operators admit two-sided inverses $G_P,G_R,G_{R^*}$ on compatible one-sided domains, with the required uniqueness and boundary data, and all the following compositions are defined there. These hypotheses involve cap and wall conditions as well as algebra.

\begin{proposition}\label{f:green-homotopy}
On that complex a contracting homotopy has three possibly nonzero components, on degrees $0,1,2$ respectively:
\[
 h_0=G_R C,\qquad h_1=G_P,\qquad h_2=C^*G_{R^*} .
\]
\end{proposition}
\begin{proof}
From $JK=0$ and $K^*J=0$ obtain $PK=C^*BR$ and $K^*P=R^*BC$. With the inverse, support and domain assumptions these identities imply $KG_R=G_PC^*B$ and $BCG_P=G_{R^*}K^*$. On $V$ one thus has $Kh_0+h_1J=G_P(C^*BC+J)=1$. On $V^\vee$ one has $Jh_1+h_2K^*=(J+C^*BC)G_P=1$. On the two remaining spaces $h_0K=G_RR=1$ and $K^*h_2=R^*G_{R^*}=1$. These are all four degrees of \eqref{f:contraction}.
\end{proof}

\begin{example}[Changing the gauge condition without changing the response class]\label{f:gauge-matrix}
Use Euclidean pairings on $U=\mathbb R$, $V=\mathbb R^2$. Set $K\lambda=(\lambda,0)$, $J=\operatorname{diag}(0,q)$, $q\ne0$. A compatible source is $(0,s)$; all its preimages under $J$ have second coordinate $s/q$. Choose $C(x,y)=a x+d y$ with $a\ne0$, and $B=b\ne0$. Then
\[
 P=\begin{pmatrix}ba^2&bad\\ bad&q+bd^2\end{pmatrix},\qquad
 \det P=ba^2q,\qquad P^{-1}(0,s)^{\mathsf T}
   =\big(-ds/(a q),\,s/q\big)^{\mathsf T}.
\]
Changing $d$ moves the first coordinate by a vector in $\operatorname{im}K$ and does not change the value seen by another closed source. Conversely a test source $(1,0)$ would read a spurious dependence on $d$. It was never closed. At $a=0$ the gauge condition fails to slice the orbit: $CK=0$ and the proposed inversion fails. At $q=0$ there is a different problem, an actual residual physical direction.
\end{example}

\section{A differential complex on spacetime}\label{f:hodge}
On oriented pseudo-Riemannian spacetime the metric gives alternating covectors a nondegenerate pairing, generally indefinite. The Hodge star $*:\Lambda^kT^*M\to\Lambda^{n-k}T^*M$ is fixed by $\alpha\wedge *\beta=\langle\alpha,\beta\rangle_g\,\epsilon_g$, where the $k$-form pairing includes the coordinate factor $1/k!$. To see existence, use an oriented orthonormal coframe: the star sends a wedge of distinct coframe covectors to the complementary wedge with the sign from the metric times the permutation sign. This description also gives $*^2=(-1)^{k(n-k)+1}$ in Lorentz signature with one minus sign. Write $d^*$ for the formal adjoint of $d$ under the integrated metric pairing on compactly supported forms without boundary. On one-forms, $d^*A=-\nabla^a A_a$ by Stokes. It is neither pointwise complex conjugation nor the field-space differential $\delta$. Adjointness and $d^2=0$ imply $(d^*)^2=0$.

Take $K=d$ on scalar parameters, $J=-d^*d$ on one-form potentials, $C=d^*$, and $B=-1$ after metric identification of dual densities. The Maxwell equation $JA=f$ requires $d^*f=0$, the conservation of its source. The gauge-fixed field and parameter operators are $P=-W_1$ and $R=W_0$ with $W_k=dd^*+d^*d$ on $k$-forms. Up to overall sign these have metric wave principal part; on globally hyperbolic spacetimes without boundary apply the wave theorem of Chapter~\ref{f:waves} and the one-sided extension on p.~\pageref{f:one-sided-extension}. For compact input supported in $K$, the retarded outputs and all their derivatives have support in $J^+(K)$, and the advanced ones in $J^-(K)$. These are the larger spaces used when a second inverse is applied. Write $S_k$ for the same-side causal inverse of $W_k$. Then $W_1 d=dW_0$ since $d^2=0$, and uniqueness with the same one-sided data gives $S_1d=dS_0$. Consequently the three homotopy blocks are $S_0d^*$, $-S_1$, and $dS_0$. On one-forms, for instance, $dS_0d^*+S_1d^*d=S_1W_1=1$. With a wall, $d$ might take a scalar in its allowed domain to a one-form outside the field domain. The algebraic identity $W_1d=dW_0$ would still hold, but the inverse intertwining conclusion would not follow from it.

Why use two one-sided homotopies? On their common compact-source input one has $D(h^R-h^A)+(h^R-h^A)D=0$. A closed source is therefore sent to a homogeneous field; changing the source by an exact one changes the field by an exact one (with the sign fixed by this identity). This is the replacement for saying that the causal propagator descends from test sources to physical solutions. If two gauge choices admit homotopies on the \emph{same} retarded complex, their two retarded responses on a closed source differ by a closed field $w$. Then $w=D\widehat h w$ is a proper gauge change there, and similarly on the advanced complex. Pairing with another closed source does not distinguish the two choices. This proof requires precisely the domain where the homotopy on $w$ and the adjoint pairing exist; a residual projector or a surface charge changes the answer.

\section*{Exercises}
\begin{exercise}In Example~\ref{f:gauge-matrix} set $q=0$. Find all the cohomology groups of the four-term complex and explain why a contraction on the whole complex is impossible.\end{exercise}
\noindent\emph{Solution.} $K$ is still injective and $K^*$ onto. At degree zero $\mathbb R^2/\mathbb R(1,0)$ is one-dimensional; at degree one $\ker K^*=\mathbb R(0,1)$ is also one-dimensional since $J=0$. Endpoint cohomologies vanish. Equation~\eqref{f:contraction} would make both middle groups vanish.
\begin{exercise}In a finite-dimensional complex with a positive inner product let $\Delta=DD^*+D^*D$. Prove that $\ker\Delta=\ker D\cap\ker D^*$. If $G$ is the inverse of $\Delta$ on the orthogonal complement of this kernel and zero on the kernel, construct a homotopy to orthogonal projection onto the kernel.\end{exercise}
\noindent\emph{Hint.} Pair $\Delta v$ with $v$. Both $D$ and $D^*$ commute with $\Delta$ on the full graded space and preserve its eigenspaces. Hence $h=D^*G$ satisfies $Dh+hD=\Delta G=1-\Pi$. No causal support condition or indefinite inner product is present in this exercise.
\begin{exercise}\label{f:gauge-boundary-ex}On a spatial region with boundary take a scalar charge density constraint $\partial_iE^i=0$. Pair a variation $\partial_i\delta E^i$ with $\lambda$ and integrate by parts. Why might a boundary value of $\lambda$ be physically measurable although the bulk transformation is a gradient?\end{exercise}
\noindent\emph{Check.} $\int\lambda\,\partial_i\delta E^i=-\int (\partial_i\lambda)\delta E^i+\int_{\partial}\lambda\,\delta E^n$. If the surface flux is variable and the parameter is allowed to be nonzero there, its pairing need not vanish. Setting the proper parameter to zero there, or imposing a flux condition, changes this test. Chapter~\ref{f:symplectic} gives its interpretation as a charge rather than a gauge degeneracy.
